\documentclass{article}
\usepackage{hyperref}
\usepackage{Style}

\nocite{*} % Comentar si quiero citar
%\addbibresource{bibliografia.bib} % Quitar el comentado si quiero usar bibliografia

\begin{document}

\begin{minipage}{2.5cm}
    \includegraphics[width=2cm]{imagen_puc.jpg}
\end{minipage}
\begin{minipage}{12cm}
    {\sc Pontificia Universidad Católica de Chile\\
    Facultad de Matemáticas\\
    Departamento de Matemática\\
    Profesor: Manuel Cabezas -- Ayudante: Benjamín Mateluna}
\end{minipage}
\vspace{2ex}

{\centerline{\bf Introducción a la Geometría - MAT1304}
\centerline{\bf Ayudantía 27 - Repaso Polinomios}}
\centerline{\bf 25 de junio de 2026}

\vspace{1cm}
\noindent\textbf{Problema 1.} Sea $p(x)$ un polinomio no constante. Mostrar que $q(x)=p(cos(x))$
no puede ser escrito como un polinomio.

\vspace{5mm}
\noindent\textbf{Problema 2.} Usar que $cos(\frac{\pi}{7})-cos(\frac{2\pi}{7})+cos(\frac{3\pi}{7})
=\frac{1}{2}$ para encontrar un polinomio en $\Q[x]$ tal que $cos(\frac{\pi}{7})$ es raíz.

\vspace{5mm}
\noindent\textbf{Problema 3.} Encontrar todas las raíces racionales de
\begin{equation*}
    p(x)=7x^{2}+\frac{7}{3}x^{4}+\frac{28}{9}x^{3}+\frac{28}{9}x^{2}-\frac{35}{9}x+\frac{7}{9}
\end{equation*}

\vspace{5mm}
\noindent\textbf{Problema 4.} Considere el polinomio $p(x)=x^{3}+x^{2}+x-3$ y demuestre que
\begin{equation*}
    \sqrt[3]{\frac{44}{27}+\sqrt{\frac{8}{3}}}+\sqrt[3]{\frac{44}{27}-\sqrt{\frac{8}{3}}}
    =\frac{4}{3}
\end{equation*}

\vspace{5mm}
\noindent\textbf{Problema 5.} Factorice completamente $x^{4}+2x^{3}+2x^{2}-2x-3$ en $\Q[x]$ y en 
$\C[x]$.

\vspace{1cm}
\noindent\textbf{\underline{Ejercicios Propuestos:}}

\vspace{5mm}
\noindent\textbf{Extra 1.} Resuelva la ecuación $z^{5}=-2+2i$.

\vspace{5mm}
\noindent\textbf{Extra 2.} Encontrar un polinomio en $\Q[x]$ tal que $\sqrt{7+\sqrt{4+\sqrt{3}}}$
es raíz.

%\printbibliography % Quitar el comentado si quiero usar bibliografia

\end{document}